\section{Implementation Stack}
\label{sec:stack}

\subsection{Model Serving}
\textbf{vLLM.} vLLM \cite{kwon2023pagedattention} is the specified serving
layer for the Linux deployment profile. Its x86 CPU backend requires Linux,
supports FP32, FP16 and BF16 with AWQ, GPTQ and INT8 W8A8 quantisation,
recommends AVX-512 and offers ``limited features'' on AVX2; pre-built CPU
wheels and the \texttt{vllm/vllm-openai-cpu} image exist. Qwen3-1.7B, 4B and
8B are on its validated CPU model list. The CPU backend implements hybrid
Gated-DeltaNet models, but the accelerated path needs AMX or AVX-512-BF16;
elsewhere it falls back to a generic path. At the time of writing the latest
published release is v0.31.0 (v0.32.0 is a tag, not a release); the
specification pins 0.31.0. vLLM exposes OpenAI-compatible chat completions
with \texttt{logprobs} and \texttt{top\_logprobs}, the \texttt{n} parameter,
a \texttt{qwen3} reasoning parser that separates \texttt{reasoning\_content},
and \texttt{chat\_template\_kwargs} for \texttt{enable\_thinking}; these are
exactly the hooks the gate needs. On the macOS Intel reference machine vLLM
runs only inside a Linux container.

\textbf{llama.cpp.} \texttt{llama-server} (release v0.7.0, October 2026) is
the mandated fallback and the primary engine on the reference machine: GGUF
weights, a mature AVX2 path, OpenAI-compatible chat completions with SSE
streaming and \texttt{logprobs}, \texttt{--reasoning-format} to separate
reasoning, \texttt{--reasoning-budget} for budget forcing, n-gram speculative
decoding without a draft model (\texttt{--spec-type ngram-simple}), and
explicit thread pinning. Both engines present the same HTTP contract, so the
Rust gateway has a single backend trait.

\textbf{Ollama} is convenient for demonstrations but its OpenAI endpoint does
not return log-probabilities, which the gate requires; it is therefore not a
supported backend.

\subsection{Fine-Tuning}
Unsloth (2026.10.6) trains with LoRA and QLoRA two to five times faster than
baseline PEFT and exports merged weights to GGUF directly. Its core trainer
requires a torch accelerator: the device-selection code raises ``You need a
GPU'' on CPU-only hosts, and the documented CPU support covers inference and
data preparation, not training. The feasible path is therefore: collect and
filter traces on the CPU machine, train on a rented GPU (a 16\,GB T4 covers
QLoRA on Qwen3-1.7B and 4B \cite{dettmers2023qlora,hu2022lora}, and BF16
LoRA on Qwen3.5 up to 4B; Unsloth advises against 4-bit training for
Qwen3.5), export GGUF, and run the promotion gate and all evaluation on CPU.
Unsloth's default hyperparameters (rank 16--32, $\alpha=r$, learning rate
$2\times10^{-4}$, one to three epochs, effective batch 16, dropout 0, all
attention and MLP projections) are the starting point of the training
specification.

\subsection{Rust Web Application}
Table~\ref{tab:crates} lists the crates and the versions verified on
crates.io on 2026-10-18. \textbf{axum} on \textbf{tokio} and \textbf{tower}
is the mainstream, well-maintained choice for an async HTTP service with
SSE (\texttt{axum::response::Sse}) and WebSockets built in;
\textbf{reqwest} with \texttt{eventsource-stream} consumes the model
servers' SSE; \textbf{serde} and \textbf{schemars} give JSON (de)serialisation
and JSON-Schema generation from the same Rust types that the OpenAPI document
is generated from with \textbf{utoipa}; \textbf{sqlx} with SQLite gives
compile-time-checked queries without a database server; \textbf{tracing}
with OpenTelemetry export and \textbf{metrics-exporter-prometheus} provide
observability; \textbf{askama} templates with \textbf{htmx} keep the
front end server-rendered and dependency-light (a single-page framework such
as Leptos is an option for a later phase, not a requirement); \textbf{goose}
provides load generation for the performance console and \textbf{criterion}
micro-benchmarks. The toolchain is pinned to the stable release current at
project start (1.96 on the reference machine; 1.99 is current upstream).

\begin{table}[t]
\caption{Rust crates (versions on crates.io, 2026-10-18)}
\label{tab:crates}
\centering\scriptsize
\begin{tabular}{@{}p{2.6cm}p{1.6cm}p{3.6cm}@{}}
\toprule
Crate & Version & Role \\
\midrule
tokio & 1.53.3 & async runtime \\
axum & 0.8.9 & HTTP routing, SSE, WebSocket \\
tower, tower-http & 0.5.3, 0.7.1 & middleware: timeouts, limits, CORS, compression \\
hyper & 1.12.0 & HTTP/1.1 and HTTP/2 \\
reqwest & 0.13.5 & client to model servers (streaming) \\
eventsource-stream & 0.2.3 & SSE parsing of upstream streams \\
serde, serde\_json & 1.0.229, 1.0.154 & (de)serialisation \\
schemars & 1.2.3 & JSON Schema from types \\
jsonschema & 0.59.0 & schema verifier \\
utoipa, utoipa-\allowbreak swagger-ui & 6.0.0, 10.0.1 & OpenAPI 3.1 generation and docs UI \\
sqlx & 0.9.0 & SQLite access, migrations \\
tracing, tracing-subscriber & 0.1.44, 0.3.23 & structured logs and spans \\
opentelemetry, tracing-\allowbreak opentelemetry & 0.33.1, 0.34.0 & trace export \\
metrics, metrics-\allowbreak exporter-\allowbreak prometheus & 0.24.6, 0.18.3 & metrics \\
askama & 0.16.1 & HTML templates \\
clap, figment & 4.6.7, 0.10.19 & CLI and layered configuration \\
thiserror, anyhow & 2.0.21, 1.0.104 & error taxonomy \\
uuid, chrono & 1.28.0, 0.4.45 & identifiers, time \\
governor & 0.10.4 & rate limiting \\
criterion, goose & 0.8.2, 0.18.1 & micro-benchmarks, load tests \\
\bottomrule
\end{tabular}
\end{table}

\subsection{Reference Machine and Measured Feasibility}
\label{sec:feasibility}
The reference machine is an Intel Core i3-8100 (four cores, no simultaneous
multithreading, AVX2 and FMA, no AVX-512), 32\,GB of DDR4, macOS~15.4.1. We
built llama.cpp from source (commit \texttt{797de63}, 2026-10-17) and ran
\texttt{llama-bench} with four threads, 512-token prefill and 128-token
decode, three repetitions. Table~\ref{tab:measured} reports the results. The
Qwen3-1.7B decode rate of $15.5$ tokens/s against a 1.26\,GiB weight file
implies an effective bandwidth of about $20$\,GB/s, roughly half the
theoretical $38.4$\,GB/s of dual-channel DDR4-2400, which is in line with the
roofline of Proposition~\ref{prop:roofline} and with published AVX2 data
points. We also confirmed, through Ollama, that Qwen3-1.7B generates in both
thinking and non-thinking modes on this machine; a 512-token thinking cap was
exhausted on a simple arithmetic question before the answer was produced,
which is the empirical reason the controller forces the answer when the
budget is reached (Proposition~\ref{prop:budget}). We also exercised \texttt{llama-server} through its OpenAI-compatible API with
Qwen3-1.7B: per-token log-probabilities are returned with \texttt{logprobs};
the \texttt{n} parameter is capped by the number of server slots, so the
gate's $k$ cheap samples are issued as $k$ parallel requests against
$k$ slots (four 12-token samples completed in 3.0\,s with 47 of 48 prompt
tokens served from the prompt cache); \texttt{enable\_thinking} toggles the
mode per request; \texttt{--reasoning-budget 256} closed the reasoning span
after 256 tokens and the model then produced an answer; and SSE streaming
delivered chunked deltas. These are the exact hooks the specification relies
on. We did not run vLLM on the
reference machine because the Docker daemon was not running; the
specification's vLLM configuration is taken from its documentation and is
marked as untested on this machine.

\begin{table}[t]
\caption{Measured throughput on the reference machine (llama.cpp, 4 threads; mean $\pm$ s.d. over 3 runs)}
\label{tab:measured}
\centering\scriptsize
\begin{tabular}{@{}llrr@{}}
\toprule
Model (GGUF) & Size & Prefill tok/s & Decode tok/s \\
\midrule
Qwen3-0.6B Q8\_0 & 604\,MiB & $180.6\pm2.4$ & $32.6\pm1.8$ \\
Qwen3-1.7B Q4\_K\_M & 1.26\,GiB & $94.1\pm9.5$ & $15.5\pm3.4$ \\
\MEASUREDROWS
\bottomrule
\end{tabular}
\end{table}

From these measurements and the roofline we derive the CPU-baseline
targets of the specification: a fast-path first token under two seconds for
prompts up to 1\,000 tokens with prompt caching, a fast-path reply of 150
tokens in about ten seconds, and a slow-path reply with a 1\,024-token
thinking budget on the 4B model in about two to three minutes; the latter
is acceptable for a demonstration and for batch evaluation, and is the
reason the roadmap keeps interactive use of \Stwo\ optional until the GPU
variant exists.
